\let\negmedspace\undefined
\let\negthickspace\undefined
\documentclass[journal]{IEEEtran}
\usepackage[a5paper, margin=10mm, onecolumn]{geometry}
%\usepackage{lmodern} % Ensure lmodern is loaded for pdflatex
\usepackage{tfrupee} % Include tfrupee package

\setlength{\headheight}{1cm} % Set the height of the header box
\setlength{\headsep}{0mm}     % Set the distance between the header box and the top of the text

\usepackage{gvv-book}
\usepackage{gvv}
\usepackage{cite}
\usepackage{amsmath,amssymb,amsfonts,amsthm}
\usepackage{algorithmic}
\usepackage{graphicx}
\usepackage{textcomp}
\usepackage{xcolor}
%\usepackage{txfonts}
\usepackage{listings}
\usepackage{enumitem}
\usepackage{mathtools}
\usepackage{gensymb}
\usepackage{comment}
\usepackage[breaklinks=true]{hyperref}
\usepackage{tkz-euclide} 
\usepackage{listings}
% \usepackage{gvv}                                        
\def\inputGnumericTable{}                                 
\usepackage[latin1]{inputenc}                                
\usepackage{color}                                            
\usepackage{array}                                            
\usepackage{longtable}                                       
\usepackage{calc}                                             
\usepackage{multirow}                                         
\usepackage{hhline}                                           
\usepackage{ifthen}                                           
\usepackage{lscape}
\usepackage{circuitikz}
\tikzstyle{block} = [rectangle, draw, fill=blue!20, 
    text width=4em, text centered, rounded corners, minimum height=3em]
\tikzstyle{sum} = [draw, fill=blue!10, circle, minimum size=1cm, node distance=1.5cm]
\tikzstyle{input} = [coordinate]
\tikzstyle{output} = [coordinate]


\begin{document}

\bibliographystyle{IEEEtran}
\vspace{3cm}

\title{5.3.24}
\author{AI25BTECH11030 -Sarvesh Tamgade}
{\let\newpage\relax\maketitle}

\renewcommand{\thefigure}{\theenumi}
\renewcommand{\thetable}{\theenumi}
\setlength{\intextsep}{10pt} 


\numberwithin{equation}{enumi}
\numberwithin{figure}{enumi}
\renewcommand{\thetable}{\theenumi}


\textbf{Question}:Solve the following pair of linear equations:
\begin{align}
   3x - 5y = 4 
\end{align}
\begin{align}
    2y + 7 = 9x
\end{align}

\textbf{Solution}: \\
The equation of a line:
\begin{align}
\vec{n^T} \vec{x} = c
\end{align}

Line \(L\):
\begin{align}
\myvec{3 & -5} \vec{x} = 4
\end{align}

Line \(K\):
\begin{align}
\myvec{9 & -2} \vec{x} = 7
\end{align}

Combining both the equations:
\begin{align}
\myvec{3 & -5 \\ 9 & -2} \vec{x} = \myvec{4 \\ 7}
\end{align}

Writing the augmented matrix:
\begin{align}
\myaugvec{2}{
3 & -5 & 4 \\
9 & -2 & 7
}
\end{align}

Perform the row operation \( R_2 \to R_2 - 3 R_1 \):
\begin{align}
\myaugvec{2}{
3 & -5 & 4 \\
0 & 13 & -5
}
\end{align}

Perform the row operation \( R_2 \to \frac{1}{13} R_2 \):
\begin{align}
\myaugvec{2}{
3 & -5 & 4 \\
0 & 1 & -\frac{5}{13}
}
\end{align}

Back substitute using \( R_1 \to R_1 + 5 R_2 \):
\begin{align}
\myaugvec{2}{
3 & 0 & \frac{27}{13} \\
0 & 1 & -\frac{5}{13}
}
\end{align}

Therefore, the solution vector is:
\begin{align}
\vec{x} = \myvec{\frac{9}{13} \\ -\frac{5}{13}}
\end{align}


\textbf{Final Answer:} Therefore, the solution vector is:
\[
\boxed{
\vec{x} = \myvec{\frac{9}{13} \\ -\frac{5}{13}}
}
\]

\begin{figure}[htbp]
    \centering
    \includegraphics[width=0.65\linewidth]{FIG/fig1.png}
    \caption{Vector Representation}
    \label{fig:FIG/fig1.png}
    \end{figure}

\end{document}  